\scanpage{143}\mspage{17}
% archive sheet number: 142
\unsure{Or|On} si $\pi^{\pi_0}=\pi$, p.ex. si $\pi=2^c$ avec $c\geq\pi_0$, on aura
\[
\deleted{\unsure{M}}\quad\mathrm{Filt}^{\pi}=M_\pi.\tag*{(3.2 bis)}
\]
(Si on admet l'hyp. du continu généralisé, on \unsure{â|a} $\pi^{\pi_0}\unsure{\ne|=|\geq}\pi$ pour \unsure{tt} $\pi>\pi_0$, donc $\mathrm{Filt}^{\pi}=M_\pi$ si $\pi\geq\pi_1$ et $\pi>\pi_0$.)

\deleted{Mon but à présent}

Je pose, pour \unsure{tt} cardinal $\pi\geq\pi_0$
\[
\mathrm{Fl}^{\pi}(M)=\{f\in\mathrm{Fl}(M)\mid\text{source et but de }f\text{ sont dans }\mathrm{Filt}^{\pi}\}\tag*{(3.3.)}
\]
-- ce qui sans doute donne une filtration cardinale sur $\uline{\mathrm{Fl}}\,M$. (mais ce n'est pas \unsure{là notre} souci\ldots). Mon but à présent, c'est de dégager des conditions sur $C_0\subset\mathrm{Fl}_{\deleted{\unsure{q|coq}}\mathrm{coquar}}(M)$, (petite partie de $\mathrm{Fl}\,M$) et sur \unsure{les|le} cardinal $\pi$, et l'ordinal $\alpha$, pour que l'on ait
\[
\begin{aligned}
f\in\mathrm{Fl}^{\pi}M\quad\overset{?}{\Longrightarrow}\quad&\deleted{\unsure{f}}X_\alpha(f)\in\mathrm{Filt}^{\pi}\\
&\text{(i.e. }p_\alpha(f)\in\mathrm{Fl}^{\pi}(M),\\
&\text{ou encore }i_\alpha(f)\in\mathrm{Fl}^{\pi}(M)\text{)}
\end{aligned}\tag*{(3.4.)}
\]
Plus généralement, je me propose de donner une \deleted{majoration} estimation de la ``cardinalité'' de $X_\alpha(f)$, pour $C_0$ \unsure{fixé|fixe}, en fonction de $C_0$, de $f$, de $\alpha$, plus exactement, en fonction des \inserted{trois} cardinaux $\pi(C_0)$, $\pi(f)$, $\pi(\alpha)$ définis comme
\[
\left\{\begin{aligned}
\pi(C_0)&=\text{plus petit cardinal tel que }C_0\subset\mathrm{Fl}^{\pi}M\\
\pi(f)&=\pi(\{f\})=\text{plus petit cardinal tel que }f\in\mathrm{Fl}^{\pi}M\\
\pi(\alpha)&=\mathrm{card}\,\alpha
\end{aligned}\right.\tag*{3.5.}
\]
Donc je cherche une fonction,
